% p1-06-multiple-testing

\section{P1 §6. Multiple Testing and Correction Procedure}

6. Multiple Testing and Correction Procedure

   Figure (fig-p1-ch6-multiple-testing). Look-elsewhere triptych --- N candidate folios / null residual
                                  histogram / Bonferroni xN penalty.

  6.1 The Multiple Testing Problem
  The symbolic search over S(C) constitutes a large-scale multiple comparison. For each
  target constant Xi $\in$ T, the minimum residual dC(Xi) is the minimum over |S(C)|
  comparisons. The minimum of M independent uniform [0,1] random variables has
  expectation 1/(M+1), decreasing as M-1. Since |S(C)| $\sim$ O(C4), residuals as small as
  O(C-4) are expected even under purely random target values.

  This is the central obstacle for symbolic constant studies: the combinatorial
  proliferation of candidates makes it statistically almost certain that small residuals will
  be found, even for randomly generated target sets. Any significance claim that does not
  explicitly account for this proliferation is methodologically void.

  6.2 Effective Trial Count
  For a grammar of complexity C, the effective number of statistical trials for each target
  constant is

  Neff(C) := |S(C)| approx (3 C4)/(2). 14

  A naively computed p-value for a single match with residual $\varepsilon$ must be inflated by the
  Bonferroni factor:

  pcorrected = min(1, Neff(C) $\cdot$ praw). 15

  At C = 10, Neff approx 15{,}000, so a raw p-value of 10-3 yields a corrected p-value of
  approximately 15, meaning the match is completely non-significant. This illustrates why
  uncorrected coincidence counting is scientifically insufficient.

  6.3 Family-Wise Error Rate Control

  For joint assessment of coverage across all N = |T| constants simultaneously, the total
  number of comparisons is N $\cdot$ Neff(C). Bonferroni correction at family-wise error rate
  (FWER) alphaFWER = 0.05 requires

  alphaindividual = (0.05)/(N $\cdot$ Neff(C)). 16

  For N = 10 constants and C = 10, this gives alphaindividual approx 3.3 $\times$ 10-7. No
  claim of individual constant significance is made unless this threshold is met.

  6.4 False Discovery Rate Control
  For the weaker goal of identifying a fraction of constants with genuine compressibility,
  the Benjamini--Hochberg (BH) procedure is applied to the ranked p-values at the
  dataset level. The BH procedure controls the expected fraction of false discoveries
  among rejections at a specified false discovery rate (FDR). It is appropriate when a
  non-trivial proportion of the targets is expected to exhibit compressibility under the
  alternative hypothesis.

  6.5 Permutation-Based Calibration
  All reported p-values are calibrated against the empirical null distribution derived from
  B $\geq$ 10{,}000 permutation resamples under Mperm. A result is reported as significant
  only if its corrected p-value is below 0.05 under both analytic Bonferroni correction and
  permutation calibration. This two-layer requirement embodies the honest-admission
  ethos of the PhD monograph's Limitations chapter (PhD: Ch.52 Limitations).

\subsection{References}

- Vasilev, D. \textit{Flos Aureus Monograph: Trinity Constants}. DOI \href{https://zenodo.org/doi/10.5281/zenodo.19227877}{10.5281/zenodo.19227877}.
- Pellis, S. "The Fine-Structure Constant and the Golden Ratio." \href{https://vixra.org/pdf/2110.0117v4.pdf}{viXra 2110.0117 v4}.
- Sherbon, M. A. "Fundamental Nature of the Fine-Structure Constant." \href{https://hal.science/hal-02341850/document}{HAL hal-02341850}.
